% A.73: a uniform finite-level obstruction, with no completed-data claim.
\begin{statement}{A.73}{Proved}{Nonconstant joint fibre characters at every exponent divisible by four}
For every integer $n$ divisible by four there is a geometrically
integral three-coordinate Kummer curve
\[
       T-c_i=(-c_i)Y_i^n\qquad(0\le i\le2)
       \tag{A.73.1}
\]
with distinct nonzero rational $c_i$, a new rational puncture
$c\notin\{0,1,c_0,c_1,c_2\}$, a rational point $P_0$ above
$T=0$, and an ordinary point $P_2\in C(\Q_2)$ above $T=1$,
having the following properties.  Write $n=2^rv$ with $v$ odd
and $r\ge2$, and put $h=2^{r-1}$.
Every field factor of the fibre at $T=c$ contains
$\Q(\mu_{2^{r+1}})$.  The even cyclotomic character
\[
 \chi_r:\operatorname{Gal}(\Q(\mu_{2^{r+1}})/\Q)
                  \longrightarrow\Q/\Z,
 \qquad \chi_r(-1)=0,\quad\chi_r(5)=\frac1h,
       \tag{A.73.2}
\]
gives a class on the smooth projective model $C$,
\[
       \mathcal A_n=
            \left(\frac{T-c}{-c},\chi_r\right)\in\operatorname{Br}C,
 \qquad \mathcal A_n(P_0)=0,
 \qquad \operatorname{inv}_2\mathcal A_n(P_2)=\frac12.
       \tag{A.73.3}
\]
The notation $(f,\chi_r)$ denotes the cyclic-algebra class, with
the convention $\operatorname{inv}_2(f,\chi_r)
=\chi_r(\operatorname{rec}_2(f))$.
In particular $\mathcal A_n$ does not belong to the image of
$\operatorname{Br}\Q\to\operatorname{Br}C$.

Set $t_2=1$ and $t_w=0$ at every other place $w$ of $\Q$.
The three global classes $-c_i$ modulo $n$ restrict to the classes
of $t_w-c_i$ at all places.  No global class modulo $h$ has
restrictions equal to the classes of $t_w-c$ at all places.
Consequently these prescribed, everywhere ordinary old data
cannot acquire the new coordinate modulo $h$, modulo $n$, or
in the completed Kummer group.
\end{statement}
\begin{proof}
The construction uses only two finite congruences and one finite
root calculation.  Put
\[
             x=1+3\cdot2^r,\qquad c=\frac1{1-x}=-\frac1{3\cdot2^r}.
       \tag{A.73.4}
\]
The element $5$ generates the subgroup of units congruent to one
modulo four in $(\Z/2^{r+2}\Z)^\times$; its order is $2^r$.
For example, the identity
$v_2(5^{2^j}-1)=j+2$ follows by induction from
$5^{2^{j+1}}-1=(5^{2^j}-1)(5^{2^j}+1)$, and proves the order
and generation by counting this subgroup.
Choose the unique $b,b'\in\{0,\ldots,2^r-1\}$ with
\[
           5^b\equiv-\frac13,
       \qquad 5^{b'}\equiv-\frac{x}{3}\pmod{2^{r+2}}.
       \tag{A.73.5}
\]
Both right-hand sides are one modulo four, so these choices exist.

We have $n\ge2^r\ge r+2$.  Hence
\[
             U=\frac{2^n-x}{3\cdot5^{b'}}
                           \in1+2^{r+2}\Z_2.
       \tag{A.73.6}
\]
The logarithm and exponential give mutually inverse group
isomorphisms $1+4\Z_2\leftrightarrow4\Z_2$.
Raising to the power $n=2^rv$ therefore gives an isomorphism
\[
 \frac{1+4\Z_2}{1+2^{n+2}\Z_2}
       \xrightarrow{\ z\mapsto z^n\ }
 \frac{1+2^{r+2}\Z_2}{1+2^{n+r+2}\Z_2}.
       \tag{A.73.7}
\]
Let $z$ be the unique integer in $[1,2^{n+2})$ congruent to one
modulo four whose $n$th power represents $U$ in the right-hand
quotient.  In particular
\[
                    x+3\cdot5^{b'}z^n
                         \equiv2^n\pmod{2^{n+r+2}}.
       \tag{A.73.8}
\]
This defines $z$ by a finite calculation for every stated $n$.

Define three nonzero rational fibre constants and the old
punctures by
\[
 D_0=5,\qquad D_1=\frac{5^b}{2^{r+n}},\qquad
 D_2=-\frac{5^{b'}z^n}{2^r},\qquad
 c_i=\frac{c}{1-D_i}\quad(0\le i\le2).
       \tag{A.73.9}
\]
The two positive numbers $D_0,D_1$ are distinct, and $D_1\ne1$:
either equality would equate a nonzero power of two to a power
of five.  The third number is negative.  Thus the $D_i$ are
distinct and none is zero or one.  The map $D\mapsto c/(1-D)$
is injective on these values, proving that the $c_i$ are distinct
and nonzero and that none equals $c$.

At $T=0$, all the root coordinates in (A.73.1) equal one.
At $T=1$, their required powers are
\[
 \begin{split}
 W_0&=x+(1-x)D_0=1-12\cdot2^r,\\
 W_1&=x+(1-x)D_1
                    =2^{-n}(2^nx-3\cdot5^b),\\
 W_2&=x+(1-x)D_2=x+3\cdot5^{b'}z^n.
 \end{split}
       \tag{A.73.10}
\]
They satisfy the exact conditions
\[
 \begin{array}{c|ccc}
 i&0&1&2\\ \hline
 v_2(W_i)&0&-n&n\\
 2^{-v_2(W_i)}W_i\pmod{2^{r+2}}&1&1&1.
 \end{array}
       \tag{A.73.11}
\]
For $W_0$ this follows from its displayed formula; for $W_1$
it follows from (A.73.5) and $n\ge r+2$; for $W_2$ it is
(A.73.8).  In particular all three numbers are nonzero, so
$1\ne c_i$.
The decomposition
\[
 \Q_2^\times=2^\Z\times\{\pm1\}\times(1+4\Z_2),\qquad
 (\Q_2^\times)^n=2^{n\Z}\times(1+2^{r+2}\Z_2)
       \tag{A.73.12}
\]
follows from the same logarithm isomorphism and the fact that
$v$ is odd.  Thus (A.73.11) constructs genuine, nonzero
$n$th roots of the $W_i$ in $\Q_2$, giving $P_2$.
Both $P_0$ and $P_2$ lie over ordinary points of the old and
enlarged punctured lines.

The fibre at $T=c$ has the constants $D_i$ in (A.73.9), since
$(c-c_i)/(-c_i)=D_i$.  Their exact joint relation is
\[
             \frac{D_2}{D_1D_0^{\,b'-b}}=-(2z)^n.
       \tag{A.73.13}
\]
In each field factor of the fibre algebra let $Z_i$ denote
the image of the root coordinate with $Z_i^n=D_i$.  Then
\[
          \theta=\frac{Z_2}{2zZ_1Z_0^{\,b'-b}}
           \quad\hbox{satisfies}\quad
          \theta^n=-1,
          \qquad (\theta^v)^{2^r}=-1.
       \tag{A.73.14}
\]
All denominators are nonzero; negative values of $b'-b$ mean
ordinary inverses in that field.  Therefore $\theta^v$ is a
primitive $2^{r+1}$st root of unity.  This proves the claimed
common fibre field by an explicit element in every factor.

For completeness, the geometric classes of $T-c_i$ are
independent modulo $n$th powers: their valuations at the three
distinct $c_i$ are the three standard basis vectors modulo $n$.
Kummer theory over $\overline\Q(T)$ thus gives a field of degree
$n^3$.  The curve (A.73.1) is geometrically integral; let $C$
be its smooth projective model.  The cyclic class in (A.73.3)
has possible residues only above $c$ and infinity.  Above $c$
the restriction of $\chi_r$ is zero by (A.73.14).
At a point $Q$ above infinity,
\[
              v_Q(T-c_i)=v_Q(T)=n v_Q(Y_i),
       \tag{A.73.15}
\]
so the pole order is divisible by $n$ and kills the character
of order $h\mid n$.  The cyclic-algebra residue formula and
purity for a smooth curve now place $\mathcal A_n$ in
$\operatorname{Br}C$
\cite[\S1.4 and Corollary 3.7.3]{AgentCTSkoro2019}.

Here is the exact character evaluation.  The group
$(\Z/2^{r+1}\Z)^\times=\langle-1\rangle\times\langle5\rangle$
has $\operatorname{ord}(5)=h$, so (A.73.2) defines a character
of exact order $h$ and conductor $2^{r+1}$.
Successive squaring, starting with $5\equiv1+4\pmod8$, gives
\[
        x\equiv1+2^r
           \equiv5^{2^{r-2}}\pmod{2^{r+1}}.
       \tag{A.73.16}
\]
The last nontrivial element of the congruence kernel has
character value $1/2$, which also verifies the stated conductor.
Under the local reciprocity map, an odd two-adic unit $u$
acts on $\mu_{2^{r+1}}$ by $\zeta\mapsto\zeta^{u^{-1}}$
\cite[I, Example 3.13]{MilneCFT}.  Since the element in
(A.73.16) has order two, its inverse is itself.  The invariant
formula for cyclic classes
\cite[III, Proposition 3.6]{MilneCFT} therefore gives
\[
       \operatorname{inv}_2(x,\chi_r)
           =\chi_r(\operatorname{rec}_2(x))
           =\frac{2^{r-2}}{2^{r-1}}=\frac12.
       \tag{A.73.17}
\]
The function $(T-c)/(-c)$ equals one at $P_0$ and $x$ at
$P_2$, proving (A.73.3).  If $\mathcal A_n$ were the pullback
of $B\in\operatorname{Br}\Q$, both evaluations at these
$\Q_2$ points would equal $\operatorname{res}_{\Q_2}(B)$.
Their different values prove the asserted failure to come
from a constant class.

The equations at $P_0$ and $P_2$ show that the global old
classes $-c_i$ have all the required local restrictions
modulo $n$.  Suppose a global class $a$ modulo $h$ supplied
the new restrictions.  After division by $-c$, its local
classes would be one at all places other than two and would
equal $x$ at two.  The global cyclic algebra
\[
                            (a/(-c),\chi_r)
                                   \in\operatorname{Br}\Q
       \tag{A.73.18}
\]
is well-defined on the class modulo $h$, because $h\chi_r=0$.
It would have invariant zero everywhere except at two, where
(A.73.17) gives $1/2$.  This contradicts the global invariant
sum \cite[Theorem 12.1.8]{AgentCTSkoro2019}.
The natural maps
\[
 \widehat{\Q^\times}\longrightarrow
       \Q^\times/(\Q^\times)^n\longrightarrow
       \Q^\times/(\Q^\times)^h
       \tag{A.73.19}
\]
commute with localization, so they give the two further
nonextension assertions.
\end{proof}

\paragraph{The first higher-character instance.}
For $n=8$ one has $r=3$, $x=25$, $b=5$, $b'=7$, and $z=445$.
Thus
\[
 \begin{split}
 c&=-\frac1{24},\qquad
 (D_0,D_1,D_2)=
      \left(5,\frac{3125}{2048},-\frac{78125\cdot445^8}{8}\right),\\
 (W_0,W_1,W_2)&=
      \left(-95,-\frac{2975}{256},
                         360405469272041107177734400\right).
 \end{split}
       \tag{A.73.20}
\]
The normalized unit of the last entry is one modulo thirty-two,
and its valuation is eight.  Relation (A.73.13) reads
$D_2/(D_1D_0^2)=-(890)^8$.
The character has order four and conductor sixteen;
the normalized new coordinate $25\equiv9\pmod{16}$ has
invariant $1/2$.  Its quadratic multiple evaluates to zero,
so testing only quadratic joint characters would miss this
specific obstruction.

\paragraph{Exact normalization and scope.}
The affine isomorphism
\[
    T'=\frac{T-c_0}{c_1-c_0},\qquad
    c_i'=\frac{c_i-c_0}{c_1-c_0},\qquad
    R'=-\frac{c_0}{c_1-c_0},\qquad
    c'=\frac{c-c_0}{c_1-c_0}
       \tag{A.73.21}
\]
sends the first two old punctures to zero and one and carries
(A.73.1) to $T'-c_i'=(R'-c_i')Y_i^n$.
It also preserves the normalized new function:
$(T'-c')/(R'-c')=(T-c)/(-c)$.
Thus the same construction applies with these two prescribed
old puncture labels.

When $40\mid n$, all squareclasses of the fibre constants
$D_i$ already have square roots in $\Q(\mu_n)$: each squareclass
is represented by a signed product of powers of two and five,
because $z^n=(z^{n/2})^2$, and
$i$, $\sqrt2=\zeta_8+\zeta_8^{-1}$, and
$\sqrt5=1+2(\zeta_5+\zeta_5^{-1})$ lie in that field.
The element (A.73.14) still supplies the additional cyclotomic
field at conductor $2^{r+1}$.
No condition on the conductors of the old coefficients $-c_i$
is asserted here.  The punctures and coefficients depend on $n$.

A.73 supplies a family of fixed finite-level local tuples.
The maps (A.73.19) prove that a completed new coordinate could
not restrict to these tuples, but the construction does not
provide completed global old coordinates with these fixed
local values.  Therefore it establishes a precise obstruction
to an unrestricted new-coordinate initialization argument;
it makes no counterexample claim for the all-rational-puncture
completed statement (G).

The finite certificate
\path{code/higher_character_obstruction.py}
checks the uniform congruence algorithm, exact relations,
local power criteria, character values, and affine maps for
the stated test exponents.  The existence for every $4\mid n$
and the infinite local roots are established by the proof.
